% Class-compatibility preamble retained for a reader excerpt of
%   Johan Hastad, Bjorn Martinsson, Tamio-Vesa Nakajima, Stanislav Zivny,
%   "A logarithmic approximation of linearly ordered colourings",
%   arXiv:2404.19556v6, CC BY 4.0.
%
% The source is a plain article; the only compatibility work here is to repeat the
% author's own declaration block mechanically and unchanged (pinned-original-source.tex
% lines 5 and 18-23: the ams packages, the plain/definition theorem styles and the
% author's four theorem environments), and to carry the two packages that the publisher
% layout and the author's cross-references need (hyperref and cleveref, the latter with
% the author's own nameinlink option). The author's font package newtx is deliberately
% not reproduced: it is a font choice for the journal layout and the excerpt is rendered
% in the delivery baseline fonts, which changes no mathematical content.
% No mathematical macro, statement, symbol or assumption is introduced, changed or
% removed by this file.
\documentclass[10pt,a4paper]{article}
\usepackage{amssymb,amsmath,amsthm}
\usepackage{hyperref}
\usepackage[nameinlink]{cleveref}

\theoremstyle{plain}
\newtheorem{theorem}{Theorem}
\newtheorem{lemma}[theorem]{Lemma}
\newtheorem{observation}[theorem]{Observation}
\theoremstyle{definition}
\title{A bounded rational solution of the homogeneous edge system yields a linearly ordered colouring with two plus log-two of the dynamic range many colours}
\author{Johan H\{\textbackslash{}aa\}stad, Bj\textbackslash{}"orn Martinsson, Tamio-Vesa Nakajima, Stanislav \textbackslash{}v\{Z\}ivn\textbackslash{}'y}
\date{arxiv:2404.19556v6; CC BY 4.0}
\begin{document}
\maketitle
\noindent\textit{Editorial note.} All delivered mathematical text is the author's own, copied from the pinned author source sources/combinatorics-fl05-a-logarithmic-approximation-of-lin/pinned-original-source.tex; no statement or proof step was rewritten, reordered or removed.\par\medskip
\noindent\textit{Editorial note.} Preamble provenance: as for the other cards of this source, the retained preamble repeats the author's ams packages, theorem styles and four newtheorem declarations (pinned source lines 5 and 18-23) unchanged and adds only hyperref and cleveref; the author's layout font package newtx is not reproduced.\par\medskip
\noindent\textit{Editorial note.} Bibliography provenance: the source has no thebibliography environment and its bibliography file bib (author source lines 354-355) is not staged anywhere in this repository, so this card deliberately counts only a statement whose statement and proof cite nothing. This card carries no bibliography, and the two citations that appear in the author's surrounding prose (to the cited preprint LNR24 and to the merger of the two independent works) are outside the counted spans and are not reproduced or reconstructed.\par\medskip
\noindent\textit{Editorial note.} The counted statement and its proof contain no citation and no cross-reference, so the card needs neither a bibliography excerpt nor a navigation stub.\par\medskip
\noindent\textit{Editorial note.} The author states just after the counted lemma that it is the same lemma used in the unbalanced case of the cited preprint; the dedup record notes this and the card still counts the author's own staged statement and proof, which is the object of this production line.\par\medskip
\noindent\textit{Editorial note.} One sentence of the counted proof reads 'The case of two vertices of odd colour is similar'; the card keeps that sentence verbatim, this being the author's own level of detail, and the reformulations are the author's own.\par\medskip
\noindent\textit{Allowed standard prerequisites (editorial).}\begin{itemize}
\item The linearly ordered (LO) colouring of a hypergraph: a colouring in which every edge has a unique maximum colour. Defined in the introduction of the paper, before source line 100; the counted proof recalls the equivalent form in which the minimum colour of each edge is unique. Fixes the property that the dyadic colour assignment in the proof has to satisfy.
\item The homogeneous system Av = 0\textasciicircum{}m attached to the edge set, and its solvability from LO 2-colourability. Set up in the section of which the counted lemma is the first statement (source lines 240-265), where the author also notes that every coordinate of the reference solution w equals -1/3 or 2/3. Supplies the solution u whose minimal and maximal absolute coordinate values M and m the colour bound 2 + log\_2(M/m) is stated in terms of.
\end{itemize}\medskip
\begin{lemma}\label{lemma:unbal}
Suppose we have a solution, $u$, to the homogeneous system where $M$
is the maximal value of the absolute value of a coordinate
and $m > 0$ is the minimal absolute value.  Then, we can
LO colour with $2+\log_2{(M/m)}$ colours.
\end{lemma}

\begin{proof} For notational convenience let us instead require that the
minimal colour in each edge should be unique. We can simply reverse
the order of the colours at the end.  By scaling we can assume $M=1$.
We use even colours for positive coordinates and we
give $x$ the colour $2\ell$ if $v_x$ is at most $2^{-(2\ell-1)}$
and strictly larger than $2^{-(2\ell+1)}$.
For negative coordinates we use $2\ell+1$ as
the colour if $v_x$ is between $-2^{-2\ell}$ (inclusive)
and $-2^{-(2\ell+2)}$ (non-inclusive).
Let us verify that this gives a correct
colouring.

Take an edge $(x,y,z)$ and suppose both $x$ and $y$
get the same colour $2\ell$.   Then by the linear equation of the edge $v_z < -2^{-2\ell}$ and thus $z$ has a colour below $2\ell$.  The case
of two vertices of odd colour is similar and as the bound on the number of colours is immediate, the lemma follows.
\end{proof}

\end{document}
